% ============================================================
% Capítulo — Trabalhos Relacionados
% ParseH2IA — Guilherme Dallman Lima, 2026
% ============================================================

\chapter{Trabalhos Relacionados}
\label{cap:relacionados}

Este capítulo revisa os trabalhos mais diretamente relacionados à presente dissertação
ao longo de três eixos: (i) sistemas de parsing de dependências para o português sob o
padrão Universal Dependencies; (ii) arquiteturas que combinam encoders BERT com mecanismos
biaffine e/ou lineares para parsing sintático; e (iii) estratégias de aprendizado
multitarefa que exploram a correlação entre tarefas morfossintáticas e sintáticas,
especialmente em línguas morfologicamente ricas. A Tabela~\ref{tab:related_works_comparison}
ao final do capítulo sintetiza as principais dimensões de comparação entre os trabalhos
revisados e a proposta desta dissertação.

% ------------------------------------------------------------
\subsection*{Towards PortParser — a Highly Accurate Parsing System for Brazilian Portuguese
             Following the Universal Dependencies Framework}
% ------------------------------------------------------------

O PortParser \cite{lopes2024towards} é o principal sistema de análise sintática de
dependências para o português brasileiro sob o padrão Universal Dependencies e constitui
o baseline central desta dissertação. O trabalho realiza uma investigação sistemática dos
fatores que influenciam o desempenho do parsing: escolha do encoder, arquitetura neural,
número de épocas de treinamento e tamanho dos embeddings.

Os autores partem de três sistemas de referência consolidados — UDPipe~1.3
\cite{straka2017tokenizing}, UDPipe~2 e Stanza \cite{qi-etal-2020-stanza} — e demonstram
que a substituição do \textit{backbone} por encoders BERT produz ganhos expressivos sobre
as arquiteturas recorrentes. A comparação entre mBERT \cite{devlin2019bert} e BERTimbau
\cite{souza2020bertimbau} mostra que o modelo monolíngue especializado no português supera
consistentemente o multilíngue, atingindo aproximadamente 99\% de acurácia em PoS tagging
e próximo de 96\% de UAS e 95\% de LAS no corpus Porttinari. O sistema vencedor combina
BERTimbau Large com uma camada BiLSTM bidirecional adicional e mecanismo de atenção
biaffine \cite{dozat2016deep} para predição conjunta de HEAD e DEPREL.

O trabalho também conduz uma análise qualitativa dos erros do parser, identificando padrões
sistemáticos associados a construções sintáticas de menor representatividade no corpus.

A contribuição desta dissertação em relação ao PortParser está em demonstrar que a camada
BiLSTM intermediária pode ser suprimida: conectar o encoder diretamente a cabeçotes de
predição com ajuste fino completo — em especial o cabeçote linear — produz resultados
competitivos (a menos de 1,1 pontos percentuais em UAS) com parametrização
significativamente mais simples.

% ------------------------------------------------------------
\subsection*{75 Languages, 1 Model: Parsing Universal Dependencies Universally}
% ------------------------------------------------------------

O trabalho de \citeauthor{kondratyuk201975} \cite{kondratyuk201975} propõe o UDify, um
único modelo treinado simultaneamente sobre 75 línguas e todos os \textit{treebanks}
disponíveis no UD~2.3, utilizando o BERT multilíngue (mBERT) como encoder compartilhado.
O UDify combina um cabeçote biaffine para HEAD e DEPREL com cabeçotes lineares para UPOS,
morfologia e lematização, em regime de \textit{hard parameter sharing} — o mesmo paradigma
de compartilhamento adotado nesta dissertação. O modelo estabeleceu o estado da arte em
parsing multilíngue na época de sua publicação e demonstrou que um único encoder pré-treinado,
sem componentes recorrentes adicionais, é suficiente para obter desempenho competitivo em
tarefas UD diversas.

Do ponto de vista metodológico, o UDify antecipa diretamente a configuração investigada
nesta dissertação: encoder pré-treinado + cabeçotes de predição especializados + MTL sobre
tarefas Universal Dependencies. As diferenças são de escopo e de foco analítico. O UDify
prioriza a cobertura de 75 línguas com um único encoder fixo (mBERT), sem comparar
arquiteturas de cabeçote diferentes; esta dissertação, por contraste, foca no português
brasileiro com quatro encoders distintos e compara sistematicamente cabeçotes linear e
biaffine em condições experimentais controladas (mesmo corpus, mesmo protocolo de busca
de hiperparâmetros). O UDify tampouco investiga a relação entre capacidade paramétrica do
cabeçote e risco de sobreajuste em corpus de dimensão limitada — questão central neste
trabalho.

Adicionalmente, o UDify inclui lematização e morfologia como tarefas adicionais, enquanto
esta dissertação restringe o MTL a UPOS, HEAD e DEPREL, as três tarefas com maior correlação
linguística direta. Por fim, o UDify não analisa a instabilidade de treinamento ao longo das
épocas por rótulo de dependência, aspecto que constitui uma das contribuições metodológicas
desta dissertação.

% ------------------------------------------------------------
\subsection*{Universal Grammatical Dependencies for Portuguese with CINTIL Data,
             LX Processing and CLARIN Support}
% ------------------------------------------------------------

O trabalho de \citeauthor{branco-etal-2022-universal} apresenta um conjunto abrangente de
recursos linguísticos para o português segundo as diretrizes UD. O estudo introduz dois
recursos principais: o CINTIL-UPos, corpus anotado com etiquetas morfossintáticas e traços
morfológicos, e o CINTIL-UDep, \textit{treebank} anotado com dependências sintáticas.
A construção dos corpora emprega dupla anotação manual com adjudicação especializada,
garantindo alta consistência linguística.

Com base nesses recursos, os autores desenvolveram o LX-UTagger, baseado em BERTimbau para
etiquetagem morfossintática, e o LX-UDParser, parser de dependências baseado em estratégias
\textit{transition-based}. Os resultados reportados atingem aproximadamente 99\% de acurácia
em PoS e cerca de 90,87\% de UAS e 88,01\% de LAS. O trabalho destaca-se pela disponibilização
pública de recursos e ferramentas voltados principalmente ao português europeu, com cobertura
também do português brasileiro.

Em relação a esta dissertação, o trabalho complementa o cenário de recursos para o português
no padrão UD e reforça que encoders baseados em BERTimbau são altamente competitivos para
tarefas morfossintáticas. Contudo, o LX-UDParser adota uma arquitetura \textit{transition-based}
— fundamentalmente diferente da abordagem \textit{graph-based} com cabeçotes lineares e
biaffine investigada aqui — e não realiza comparação entre arquiteturas de cabeçote nem
experimentos de aprendizado multitarefa.

% ------------------------------------------------------------
\subsection*{Etiquetagem Morfossintática Multigênero para o Português do Brasil Segundo
             o Modelo Universal Dependencies}
% ------------------------------------------------------------

O estudo de \citeauthor{stilpetrogold} investiga a etiquetagem morfossintática multigênero
para o português brasileiro usando o padrão UD, com foco em domínios textuais além do
jornalístico: textos acadêmicos e conteúdo gerado por usuários. Os autores comparam
arquiteturas BiLSTM com modelos pré-treinados, demonstrando que o BERTimbau apresenta o
melhor desempenho, alcançando aproximadamente 99\% de acurácia em PoS tagging. O estudo
também evidencia que modelos treinados em múltiplos gêneros apresentam maior robustez em
cenários de transferência entre domínios.

A relevância deste trabalho para a presente dissertação reside em dois aspectos. Primeiro,
ele confirma empiricamente a superioridade das representações contextualizadas do BERT sobre
arquiteturas recorrentes para tarefas morfossintáticas em português, respaldando a escolha
dos encoders avaliados no ParseH2IA. Segundo, os resultados sobre generalização multigênero
fornecem contexto sobre os limites do corpus Porttinari-base — majoritariamente jornalístico
— e apontam para o risco de sobreajuste ao domínio de treinamento, o que reforça a
importância da regularização implícita provida pelo cabeçote linear.

% ------------------------------------------------------------
\subsection*{Multi-task Learning by Using Contextualized Word Representations
             for Syntactic Parsing of a Morphologically Rich Language}
% ------------------------------------------------------------

O trabalho de \citeauthor{ehsan2025multi} propõe uma arquitetura multitarefa para parsing
sintático aplicada ao urdu, língua morfologicamente rica e com ordem SOV. O estudo combina
parsing de constituintes e parsing de dependências em um único modelo baseado em BiLSTM,
explorando representações compartilhadas entre estruturas sintáticas distintas por meio de
embeddings contextuais ELMo. Os autores investigam tanto o regime de tarefa única quanto o
multitarefa, demonstrando que o compartilhamento de representações sintáticas melhora
consistentemente os resultados em LAS e reduz ambiguidades entre relações sintáticas
semanticamente próximas.

O trabalho fornece evidência empírica importante de que estratégias multitarefa são eficazes
para parsing em línguas morfologicamente ricas, embasando a escolha do MTL nesta dissertação.
A análise das matrizes de erro e das ambiguidades sintáticas também é metodologicamente afim
à análise de instabilidade por rótulo conduzida aqui. A principal diferença reside na
tecnologia empregada: \citeauthor{ehsan2025multi} utilizam BiLSTM e ELMo — arquiteturas
anteriores à consolidação dos Transformers — enquanto esta dissertação baseia-se
inteiramente em encoders BERT com ajuste fino completo, o que elimina a necessidade de
componentes recorrentes separados e permite comparação direta entre cabeçotes de predição.

% ------------------------------------------------------------
\begin{table*}[htbp]
\centering
\caption{Comparação entre os trabalhos relacionados e a proposta desta dissertação.}
\label{tab:related_works_comparison}
\renewcommand{\arraystretch}{1.5}
\setlength{\tabcolsep}{4pt}

\resizebox{\textwidth}{!}{%
\begin{tabular}{p{3.0cm} p{3.2cm} p{1.6cm} p{1.0cm} p{0.8cm} p{4.8cm}}
\toprule
\textbf{Trabalho} &
\textbf{Arquitetura} &
\textbf{Língua} &
\textbf{MTL} &
\textbf{UD} &
\textbf{Relação com esta dissertação} \\
\midrule

PortParser \cite{lopes2024towards} &
BERTimbau Large + BiLSTM + Biaffine &
PT-BR &
Parcial &
Sim &
Baseline principal. Referência de desempenho para parsing UD no português brasileiro.
Esta dissertação suprime o BiLSTM e compara cabeçotes linear e biaffine. \\

UDify \cite{kondratyuk201975} &
mBERT + Biaffine (HEAD/DEPREL) + Linear (UPOS/morfo) &
75 línguas &
Sim &
Sim &
Pioneiro no paradigma encoder BERT + MTL + biaffine para UD. Não compara linear vs.\
biaffine nem analisa efeito do tamanho do corpus sobre a escolha de cabeçote. \\

CINTIL / LX-UDParser \cite{branco-etal-2022-universal} &
BERTimbau + \textit{Transition-based} &
PT-EU/BR &
Não &
Sim &
Referência em recursos linguísticos UD para o português.
Arquitetura \textit{transition-based} contrasta com a abordagem \textit{graph-based} desta dissertação. \\

Stil/PetroGold \cite{stilpetrogold} &
BiLSTM / BERTimbau &
PT-BR &
Não &
Parcial &
Evidencia a importância de representações contextualizadas e a robustez multigênero
para tarefas morfossintáticas. \\

\citeauthor{ehsan2025multi} \cite{ehsan2025multi} &
BiLSTM + ELMo &
Urdu &
Sim &
Sim &
Demonstra empiricamente os benefícios do MTL para parsing em línguas morfologicamente ricas.
Usa arquitetura pré-Transformer. \\

\midrule

\textbf{Esta dissertação} &
BERT\textsuperscript{*} + Linear / Biaffine &
PT-BR &
Sim &
Sim &
Suprime o BiLSTM do PortParser; compara 4 encoders $\times$ 2 cabeçotes $\times$ 2 regimes
de treinamento (MTL e ablação sem UPOS); propõe análise de instabilidade por rótulo ao
longo das épocas. \\

\bottomrule
\multicolumn{6}{l}{\footnotesize \textsuperscript{*}BERTimbau Large, BERTimbau Base, mBERT e JabuticaBERT.}
\end{tabular}%
}
\end{table*}
